% §11 Problem 1: simple pendulum, swinging between the amplitudes -phi_0 and phi_0
\documentclass{SolutionFigure}

\begin{document}
\begin{tikzpicture}
    % Parameters
    \def\l{3}           % length l
    \def\phiA{28}       % angle phi
    \def\phiM{60}       % amplitude phi_0

    \coordinate (O) at (0,0);
    \coordinate (m) at ($(O) + (-90+\phiA:\l)$);
    \coordinate (m0) at ($(O) + (-90+\phiM:\l)$);
    \coordinate (v) at (0,-1);

    % Vertical, swing between -phi_0 and phi_0, and turning point
    \draw[guide] (O) -- (0,-\l-0.4);
    \draw[guide] (-90-\phiM:\l) arc[start angle=-90-\phiM, end angle=-90+\phiM, radius=\l];
    \draw[guide] (O) -- (m0);
    \pic at (-2.2,-0.2) {gravity};

    % Angles
    \pic[angle arc, angle radius=1cm, angle eccentricity=1.4, "$\phi$"] {angle = v--O--m};
    % label of phi_0 between the pendulum and the turning point, clear of the rod
    \pic[angle arc, angle radius=2cm] {angle = v--O--m0};
    \node at ({-90+(\phiA+\phiM)/2}:2.35) {$\phi_0$};

    % Pendulum
    \draw[rod] (O) -- node[pos=0.5, right] {$l$} (m);
    \node[pivot] at (O) {};
    \node[mass, label=below:$m$] at (m) {};
\end{tikzpicture}
\end{document}
